\section{Adaptive terminology memory}
\label{sec:adaptive-memory}

Separating an offline catalog from a small working set lets AutoTerm-SST
cover broad domains while keeping per-chunk search bounded. Each topic slice
packages about 1,000 term--translation pairs, a prebuilt MaxSim index over the
source terms, and a prototype: a short English topic description with example terms and
its Chinese translation. Because each slice stores its own translations, the same source term can
map to different targets in different slices.

\paragraph{Offline catalog.}
Adding a domain touches only the catalog manifest. An operator registers a
slice with one entry holding its label, term count, prototype, and index path. The router reads only the manifest and
prototypes, and the retriever skips every index outside the working set.

\paragraph{Language-model routing.}
The router asks a small instruction-tuned LLM which slice matches the talk at
the current point. Because the speech LLM translates audio directly and
produces no source transcript, routing reads target-language output, which
becomes more topic-specific as context builds. At chunk $t$, the router
concatenates the last $w=4$ translated chunks into a context $c_t$ and builds a
multiple-choice prompt that lists the prototype of every slice under a letter.
The prototype $P_d$ of slice $d$ is a one-line description followed by an
English sentence naming the slice's core topics and representative terms. The
score of slice $d$ is the probability that the router LLM $\pi$, the 4-billion-parameter Qwen3-4B-Instruct-2507~\citep{yang2025qwen3}, assigns to the letter of
$d$ as the first answer token. Averaging over the catalog order $\sigma$ and its reverse $\bar\sigma$ places
every slice at mirrored positions and offsets the preference for early or late
positions:
\begin{equation}
  o_t(d)=\tfrac12\bigl[\pi(\ell_\sigma(d)\mid c_t)
  +\pi(\ell_{\bar\sigma}(d)\mid c_t)\bigr],
\end{equation}
where $\ell_\sigma(d)$ is the letter of $d$ under order $\sigma$. Each update
takes two forward passes and needs neither generation nor training.

\paragraph{Topic filtering.}
A sticky hidden Markov model accumulates these noisy scores into a stable
posterior $b_t$ over the $N$ catalog slices. Between updates, the topic stays with
probability $1-\rho$ and otherwise moves to one of the other $N-1$ slices
uniformly, and each update treats $o_t$ as the observation likelihood:
\begin{equation}
\begin{aligned}
  b_t(d)\propto{}&\Bigl[(1-\rho)\,b_{t-1}(d)\\
  &+\tfrac{\rho}{N-1}\bigl(1-b_{t-1}(d)\bigr)\Bigr]\,o_t(d)^{\beta_t}.
\end{aligned}
\end{equation}
Consecutive contexts can share chunks, so the exponent
$\beta_t=\min(1,n_t/w)$, with $n_t$ the number of chunks translated since the
previous update, weights each update by the share of its context that is new. The switch rate $\rho$ is
the filter's only free parameter. We set $\rho=10^{-3}$, and
Appendix~\ref{app:router} shows that top-four coverage moves by at most 0.5
points as $\rho$ ranges from $10^{-5}$ to $10^{-1}$. Routing is causal, so a
decision made at chunk $t$ affects only later chunks.

\paragraph{Working-set update.}
The working set holds the most probable slices and lets outgoing and incoming
topics coexist during a transition. Every five seconds, the router scores every
catalog slice, updates the posterior, and fills a working set of at most $K=4$
slices and $B=4{,}000$ searchable terms in order of posterior probability.
Lower-ranked slices are excluded from retrieval, and a stale slice leaves once
it drops in rank. Appendix~\ref{app:config} lists all router settings.

\paragraph{Fixed evidence budget.}
All active slices share one retrieval and prompt budget. The retriever splits
100 candidate slots equally among the active slices, removes exact duplicate
pairs, reranks the rest globally, and admits at most ten into the prompt, fewer
when fewer candidates clear the 0.78 threshold. Gains of AutoTerm-SST therefore
come from placing more relevant mappings inside the same budget.

\paragraph{Non-blocking updates.}
A background worker keeps routing off the decoding path. In one batch, it scores every session whose context has changed. Each chunk
retrieves from the working set that is current when its decoding begins, so an
update takes effect at the first chunk that starts after the update finishes. Indexes of likely slices are likewise
preloaded in the background.
